
\SGASection{1}{Lifting finite subgroups}
\SGARefTarget{sga3:xv:section:1}

\SGASubsection{1.1}{Finite, smooth, central subgroups of multiplicative type}
\SGARefTarget{sga3:xv:subsection:1:1}

\SGAStatement{Proposition 1.1}{Let \(S\) be an affine scheme, let
\(S_0\) be a closed subscheme of \(S\) defined by a square-zero ideal,
let \(G\) be an \(S\)-group prescheme, and let \(H_0\) be a subgroup
scheme of \(G_0=G\times_S S_0\) that is smooth and of finite type over
\(S_0\), and of multiplicative type. A subgroup scheme of \(G\), of
multiplicative type, lifting \(H_0\) exists if and only if there exists
a subscheme \(H'\) of \(G\), flat over \(S\), lifting \(H_0\).}
\SGARefTarget{sga3:xv:proposition:1:1}

The necessity of the condition is clear; let us prove its sufficiency.
The group \(H_0\) of multiplicative type is quasi-isotrivial (\SGARefLink{sga3:x:corollary:4-5}{X~4.5}).
By Exposé \SGARefLink{sga3:x:proposition:2:1}{X~2.1}, there exist an \(S\)-group \(H\) of multiplicative
type and an \(S_0\)-isomorphism of groups
\[
  u_0:H\times_S S_0\xrightarrow{\sim}H_0.
\]
Since \(H'\) is flat over \(S\) and \(H'\times_S S_0\) is of finite
presentation over \(S_0\) (Exposé \SGARefLink{sga3:ix:proposition:2:1:item:b:01}{IX~2.1(b)}), \(H'\) is of finite
presentation over \(S\). Its fibres are smooth, so \(H'\) is smooth
over \(S\) (EGA IV~17.5.1). Since \(S\) is affine, \(u_0\) therefore
lifts to an \(S\)-morphism of preschemes
\[
  u:H\longrightarrow H'.
\]
It then follows from Exposé \SGARefLink{sga3:iii:theorem:2-1}{III~2.1} and Exposé \SGARefLink{sga3:ix:theorem:3:1}{IX~3.1} that the composite
\[
  v_0:H\times_S S_0\xrightarrow{\sim}H_0\longrightarrow G_0
\]
also lifts to an \(S\)-homomorphism
\[
  v:H\longrightarrow G.
\]
Since \(v_0\) is an immersion, so is \(v\). The image of \(H\) under
\(v\) is therefore a subgroup scheme of \(G\), of multiplicative type,
lifting \(H_0\).

\SGAStatement{Proposition 1.2}{Let \(S\) be a prescheme, let \(S_0\)
be a subprescheme of \(S\) defined by a locally nilpotent sheaf of
ideals, let \(G\) be an \(S\)-group prescheme flat and of finite
presentation over \(S\), and let \(H_0\) be a subgroup scheme of
\(G_0=G\times_S S_0\) that is smooth, finite over \(S_0\), of
multiplicative type, and central. Then there exists a unique subgroup
scheme \(H\) of \(G\), of multiplicative type, lifting \(H_0\).
Moreover, \(H\) is central. (See XVII, Appendix III, 1.)}
\SGARefTarget{sga3:xv:proposition:1:2}

\SGAInlineRefTarget{sga3:xv:proposition:1-2-bis}\SGAStatement{Proposition 1.2 bis}{Let \(S,G,S_0,G_0\) be as above,
let \(H\) be an \(S\)-group scheme of multiplicative type, smooth and
finite over \(S\), and let
\(u_0:H_0=H\times_S S_0\to G_0\) be a central homomorphism. Then
\(u_0\) lifts uniquely to a homomorphism \(u:H\to G\). Moreover,
\(u\) is central.}
\SGARefTarget{sga3:xv:proposition:1:2:bis}

The existence of the lifting \(u\) in \SGARefLink{sga3:xv:proposition:1:2:bis}{Proposition 1.2 bis} follows
readily from \SGARefLink{sga3:xv:proposition:1:2}{Proposition 1.2} by considering the graph of \(u_0\). The
lifting \(u\) is unique and central by Exposé \SGARefLink{sga3:ix:corollary:3:4}{IX~3.4} and Exposé \SGARefLink{sga3:ix:theorem:5:1}{IX~5.1}.

\SGAProof{Proof of \SGARefLink{sga3:xv:proposition:1:2}{Proposition 1.2}}
The uniqueness of \(H\), and the fact that \(H\) is central, follow
from Exposé \SGARefLink{sga3:ix:theorem:5:6:item:b}{IX~5.6(b)} and Exposé \SGARefLink{sga3:ix:corollary:3:4:bis}{IX~3.4 bis}. In view of uniqueness, to
prove existence we may assume that \(S\) is affine and that \(S_0\) is
defined by a square-zero ideal; by \SGARefLink{sga3:xv:proposition:1:1}{Proposition 1.1} it is enough to find
a subscheme of \(G\), flat over \(S\), lifting \(H_0\).

Since \(H_0\) is smooth and finite over \(S_0\), after possibly
restricting \(S\) we may assume that there is an integer \(n>0\),
invertible on \(S\), such that \(H_0={}_nH_0\). Consider the \(n\)-th
power morphism on \(G\):
\[
  u:G\longrightarrow G,\qquad x\longmapsto x^n.
\]
We continue to denote by \({}_nG\) the ``kernel of \(u\)'', that is,
the inverse image under \(u\) of the identity section of \(G\).
Notice that \(u\) is not, in general, a group homomorphism. We first
admit the following lemma.

\SGAStatement{Lemma 1.3}{Let \(k\) be a field, let \(G\) be a group
scheme locally of finite type over \(k\), let \(n\) be an integer prime
to the characteristic of \(k\), and let \(u\) be the \(n\)-th power
morphism on \(G\). Then \(u\) is étale at every point \(x\) of \(G\)
that belongs to the centre of \(G\).}
\SGARefTarget{sga3:xv:lemma:1:3}

Since \(G\) is flat and of finite presentation over \(S\), the
preceding lemma and EGA IV~17.8.2 show that, if \(x\in G\) lies above
\(s\in S\) and belongs to the centre of \(G_s\), then \(u\) is étale
at \(x\). If in addition \(x\in {}_nG\), then \({}_nG\) is étale over
\(S\) at \(x\). By hypothesis, \(H_0\) is central and is contained in
\({}_nG_0\); it is therefore contained in the largest open subset
\(V\) of \({}_nG\) that is étale over \(S\). Since \(H_0\) and
\(V\times_S S_0=V_0\) are étale over \(S_0\), \(H_0\) is an open
subscheme of \(V_0\) (EGA IV~17.8.7 and 17.9.1). The open subprescheme
of \(V\) having the same underlying space as \(H_0\) is consequently
a subscheme of \(G\), flat over \(S\), lifting \(H_0\).

It remains to prove \SGARefLink{sga3:xv:lemma:1:3}{Lemma 1.3}. The largest open subset of \(G\) on
which \(u\) is étale is invariant under extension of the base field
(EGA IV~17.8.2), so we may reduce to the case in which \(x\) is
\(k\)-rational. Let \(t_x\) denote translation by \(x\), a
\(k\)-automorphism of \(G\). Since \(x\) belongs to the centre of
\(G\), the diagram
\[
\begin{tikzcd}[column sep=large,row sep=large]
G \arrow[r,"u"] \arrow[d,"t_x"'] & G \arrow[d,"t_x^n"] \\
G \arrow[r,"u"] & G .
\end{tikzcd}
\]
is commutative. It is therefore enough to prove that \(u\) is étale
at the identity; this was established in VII\textsubscript{A}~\SGARefLink{sga3:xv:definition:8-4}{8.4}.
